%! The Tangent Function — Unit 7, Lesson 7.1 — Group Activity
\documentclass[10pt]{article}
\usepackage{precalculus-article}
\usepackage{precalculus-boxes}
\newcommand{\TallMath}[1]{$\displaystyle #1\rule[-1.4em]{0pt}{3.2em}$}

\begin{document}

\pageheader{Unit 7, Lesson 7.1}{Group Activity}
\namepartnerperiod

\vspace{0.08in}

\begin{scenariobox}[The Tangent Function --- Tiered Practice]{sky}
The tangent function connects the unit circle to slope and to periodic asymptotic
behavior.  Choose \textbf{one tier} below based on your readiness level.
Everyone should be prepared to explain their reasoning to the class.
\end{scenariobox}

\vspace{0.06in}

%----------------------------------------------------------------------
% Tier R
%----------------------------------------------------------------------
\begin{tcolorbox}[colback=white, colframe=black!40,
    title=\textbf{Tier R --- Remediate},
    fonttitle=\bfseries, arc=2mm, left=3mm, right=3mm, top=2mm, bottom=2mm]

\textbf{Context:} The unit circle values below are given.  Use them to compute
$\tan\theta = \sin\theta / \cos\theta$ at each angle.

\vspace{6pt}
\renewcommand{\arraystretch}{1.55}
\begin{tabular}{|c|c|c|c|}
\hline
$\theta$ & $\sin\theta$ & $\cos\theta$ & $\tan\theta$ \\
\hline
$0$          & $0$          & $1$          & \blank{2cm} \\
\hline
$\pi/6$      & $1/2$        & $\sqrt{3}/2$ & \blank{2cm} \\
\hline
$\pi/4$      & $\sqrt{2}/2$ & $\sqrt{2}/2$ & \blank{2cm} \\
\hline
$\pi/3$      & $\sqrt{3}/2$ & $1/2$        & \blank{2cm} \\
\hline
$\pi/2$      & $1$          & $0$          & \blank{2cm} \\
\hline
$3\pi/4$     & $\sqrt{2}/2$ & $-\sqrt{2}/2$& \blank{2cm} \\
\hline
$\pi$        & $0$          & $-1$         & \blank{2cm} \\
\hline
\end{tabular}

\vspace{8pt}
\textbf{1.}\quad At which angles above is $\tan\theta$ \emph{undefined}?  Why?
\writelines{2}

\vspace{4pt}
\textbf{2.}\quad The tangent function has vertical asymptotes at $\theta = \pi/2 + k\pi$
for integer $k$.  List the three asymptotes closest to $\theta = 0$ (one negative,
one positive, one further positive).
\writelines{2}

\vspace{4pt}
\textbf{3.}\quad Between $\theta = 0$ and $\theta = \pi/2$, is $\tan\theta$
increasing or decreasing?  How do you know from the table?
\writelines{2}

\end{tcolorbox}

\vspace{0.06in}

%----------------------------------------------------------------------
% Tier A
%----------------------------------------------------------------------
\begin{tcolorbox}[colback=white, colframe=black!40,
    title=\textbf{Tier A --- Approaching Proficiency},
    fonttitle=\bfseries, arc=2mm, left=3mm, right=3mm, top=2mm, bottom=2mm]

\textbf{Context:} The graph below shows one and a half periods of a tangent function
$f(\theta) = a\tan(\theta) + d$.

\vspace{6pt}
\begin{center}
\begin{tikzpicture}[xscale=1.6, yscale=0.45]
  \draw[->] (-1.8,0) -- (4.2,0) node[right]{\small $\theta$};
  \draw[->] (0,-4.5) -- (0,6.5) node[above]{\small $y$};
  % asymptotes
  \draw[dashed,slate] (-1.571,-4.3)--(-1.571,6.3);
  \draw[dashed,slate] ( 1.571,-4.3)--( 1.571,6.3);
  \draw[dashed,slate] ( 3.14+1.571*0,-4.3)--( 3.14+1.571*0,6.3);
  % branches: g = 2*tan(theta) + 1
  \draw[navy,thick,domain=-1.45:-0.01,samples=80,smooth]
        plot (\x,{2*tan(deg(\x))+1});
  \draw[navy,thick,domain=0.01:1.45,samples=80,smooth]
        plot (\x,{2*tan(deg(\x))+1});
  \draw[navy,thick,domain=1.70:3.10,samples=80,smooth]
        plot (\x,{2*tan(deg(\x))+1});
  % key points annotated
  \fill[navy] (0,1) circle (2pt) node[left,font=\scriptsize]{$(0,1)$};
  \fill[navy] (0.785,3) circle (2pt) node[right,font=\scriptsize]{$(\pi/4,3)$};
  % y-axis labels
  \foreach \y in {-4,-2,2,4,6}
    \draw (0.04,\y)--(-0.04,\y) node[left,font=\scriptsize]{$\y$};
  \foreach \x/\l in {-1.571/$-\frac{\pi}{2}$,0.785/$\frac{\pi}{4}$,1.571/$\frac{\pi}{2}$,3.142/$\pi$}
    \draw (\x,0.1)--(\x,-0.1) node[below,font=\scriptsize]{\l};
\end{tikzpicture}
\end{center}

\vspace{4pt}
\textbf{1.}\quad What is the period of $f$? \writeline

\vspace{4pt}
\textbf{2.}\quad List all vertical asymptotes visible in the graph.
\writeline

\vspace{4pt}
\textbf{3.}\quad On the interval $(-\pi/2, \pi/2)$, is $f$ concave up or concave down
for $\theta < 0$?  For $\theta > 0$?
\writelines{2}

\vspace{4pt}
\textbf{4.}\quad The function passes through $(0, 1)$ and $(\pi/4, 3)$.  Use these
points to find the value of $a$ and $d$ in $f(\theta) = a\tan(\theta) + d$.
Show your work.
\writelines{3}

\end{tcolorbox}

\vspace{0.06in}

%----------------------------------------------------------------------
% Tier E
%----------------------------------------------------------------------
\begin{tcolorbox}[colback=white, colframe=black!40,
    title=\textbf{Tier E --- Extension},
    fonttitle=\bfseries, arc=2mm, left=3mm, right=3mm, top=2mm, bottom=2mm]

\textbf{Context:} Analyze $g(\theta) = 2\tan\!\left(\dfrac{\theta}{2} - \dfrac{\pi}{4}\right) + 1$.

\vspace{6pt}
\textbf{1.}\quad Rewrite $g$ in the form $a\tan(b(\theta + c)) + d$ by factoring
out $b = 1/2$ from the argument.  Identify $a$, $b$, $c$, and $d$.
\writelines{3}

\vspace{4pt}
\textbf{2.}\quad State the period of $g$.  Show the computation.
\writelines{2}

\vspace{4pt}
\textbf{3.}\quad Find all vertical asymptotes of $g$.
(Set $\tfrac{1}{2}(\theta - \tfrac{\pi}{2}) = \tfrac{\pi}{2} + k\pi$ and solve for $\theta$.)
\writelines{3}

\vspace{4pt}
\textbf{4.}\quad Compute $g(0)$ and $g\!\left(\dfrac{\pi}{2}\right)$.  Show your work.
\writelines{3}

\vspace{4pt}
\textbf{5.}\quad Describe in words how the graph of $g$ is related to the parent
graph $y = \tan\theta$, addressing each transformation parameter.
\writelines{3}

\end{tcolorbox}

\end{document}
